% !TeX root = ../tesis.tex

El objetivo de este capítulo es aislar y formalizar la capa de certificación sobre la que se apoya la arquitectura del \zkBridge. La decisión que luego guiará el capítulo de diseño es la siguiente: en lugar de demostrar dentro de un circuito \ZK la validez completa del consenso nativo de \Cardano, usamos \Mithril como mecanismo intermedio de autenticación de estado \cite{cryptoeprint:2021/916,MithrilDocsOverview}.

La notación que usamos en este capítulo es una abstracción propia de la tesis. Esta formalización constituye además un aporte adicional del trabajo, porque separa la semántica criptográfica relevante de \Mithril de los detalles accidentales de su implementación. Su propósito es la semántica criptográfica \textbf{que necesita el \zkBridge}: participantes ponderados por \textit{stake}, registros autenticados, mensajes certificados, firmas agregadas, testigos, entradas públicas y relaciones verificables por circuitos. En particular, \textbf{no} vamos a reproducir uno a uno los tipos, módulos o nombres internos de la implementación de \Mithril. Cuando la implementación introduce detalles de serialización, \textit{dispatch} por era, compatibilidad de hashes o estructuras auxiliares del \MithrilAggregator, los mencionamos sólo en la medida en que afectan esa semántica.

Como motivación, observamos que por un lado, la validación directa de \Praos traería al \zkBridge detalles de alto costo y fuerte inestabilidad semántica, como \VRF, firmas \KES, nonces de época, selección de líderes y reglas de cadena. Por otro lado, \Mithril expone una interfaz más adecuada para verificación compacta: registros cerrados de participantes, \textit{aggregate verification keys} (\AVK), cadenas de certificados, y mensajes de protocolo con estructura fija. Usaremos a \Mithril como una abstracción criptográfica del consenso de \Cardano suficientemente fuerte para autenticar el estado relevante y suficientemente compacta para ser reutilizada por la capa de aplicación del \zkBridge.

\Mithril no sustituye al consenso de \Cardano en sentido pleno: no decide qué cadena es válida ni produce bloques. Hace el trabajo de un \LightClient: certificar ciertos datos derivados del estado de \Cardano, de manera que un sistema externo pueda tratarlos como objetos autenticados. Los artefactos relevantes para este proyecto son los certificados \StakeDistribution y \TransactionSnapshot \cite{MithrilDocsCertChain,MithrilDocsClient}.

La capa de actualización mantiene en \ChainB un \UTxO canónico de \StakeDistribution, inicializado con un \GenesisCertificate y luego renovado mediante certificados estándar de \Mithril. En el flujo final, las actualizaciones estándar quedan ligadas a un \ProofReceipt producido por la verificación \STM. La capa de aplicación consume como evidencia un \TransactionSnapshotCertificate autenticado por el dominio \texttt{cardano\_transactions} y una prueba de inclusión de la \LockTx correspondiente en dicho \TransactionSnapshotCertificate.

El capítulo dedicado a \Cardano motivó la necesidad de una representación verificable del estado remoto sin reejecutar todo \Praos. Lo que falta justificar ahora es por qué el \UTxO de \StakeDistribution puede considerarse auténtico, cómo se encadena con el estado anterior del \UTxO, qué información compromete la \AVK, qué significa que un certificado de \Mithril esté correctamente formado y de qué manera dicha información llega hasta los argumentos \ZK verificados.

No vamos a hablar sobre \textbf{todo} el protocolo \Mithril. Primero describimos el protocolo \STM y sus esquemas de firma, luego formalizamos la cadena de certificados, vemos qué certificados son relevantes al \zkBridge y la interfaz expuesta por el \MithrilAggregator de \Mithril. Sobre esa base, formulamos la relación $\NP$ del circuito \STM y vemos su implementación en \HaloTwo con \KZG. Una vez fijado todo esto, el capítulo \emph{Decisiones de Diseño} puede decidir qué estado y qué evidencia debe usar el \zkBridge, y el capítulo de implementación puede mostrar cómo esas decisiones se concretan con \ProofReceipt de verificación \STM y argumentos \Groth auxiliares.
